% =============================================================================
% main.tex — Hauptdokument der IDPA-Dokumentation
% Geotracker-Projekt, IDPA 2026
% =============================================================================

\documentclass[
  12pt,           % Schriftgrösse 12 pt (zulässig: 10–12 pt)
  a4paper,        % Papierformat A4
  oneside,        % Einseitiger Druck
  bibliography=totoc % Literaturverzeichnis ins Inhaltsverzeichnis
]{scrreprt}

% -----------------------------------------------------------------------------
% Sprache & Zeichenkodierung
% -----------------------------------------------------------------------------
\usepackage[T1]{fontenc}          % T1-Zeichenkodierung für Westeuropa
\usepackage[utf8]{inputenc}       % UTF-8-Eingabe
\usepackage[ngerman]{babel}       % Deutsche Sprachregeln (neue Rechtschreibung)

% -----------------------------------------------------------------------------
% Schrift & Layout
% -----------------------------------------------------------------------------
\usepackage{lmodern}              % Skalierbare Latin Modern-Schriften
\usepackage{microtype}            % Mikrotypografie (besserer Blocksatz)
\usepackage[
  left=3cm,
  right=2.5cm,
  top=2.5cm,
  bottom=2.5cm
]{geometry}                       % Seitenränder

% Zeilenabstand 1.2 (entspricht 1.2-fachem Abstand, Vorgabe 1–1.2)
\usepackage{setspace}
\setstretch{1.2}

\usepackage{parskip}              % Absatzabstand statt Einrückung

% -----------------------------------------------------------------------------
% Grafiken & Farben
% -----------------------------------------------------------------------------
\usepackage{graphicx}             % Grafiken einbinden (\includegraphics)
\usepackage[dvipsnames]{xcolor}   % Farben
\graphicspath{{./images/}}        % Standardpfad für Bilder

% -----------------------------------------------------------------------------
% Tabellen & Listen
% -----------------------------------------------------------------------------
\usepackage{booktabs}             % Professionelle Tabellen
\usepackage{longtable}            % Mehrseitige Tabellen
\usepackage{enumitem}             % Angepasste Listen

% -----------------------------------------------------------------------------
% Code-Listings
% -----------------------------------------------------------------------------
\usepackage{listings}             % Quellcode-Listings
\usepackage{lstautogobble}        % Einrückung automatisch entfernen
\lstset{
  basicstyle=\ttfamily\small,
  breaklines=true,
  frame=single,
  autogobble=true,
  language=C,
  commentstyle=\color{gray},
  keywordstyle=\color{blue},
  stringstyle=\color{red!60!black},
}

% -----------------------------------------------------------------------------
% Verweise & Links
% -----------------------------------------------------------------------------
\usepackage[
  hidelinks,       % Keine farbigen Rahmen um Links
  pdftitle={IDPA - Geotracker},
  pdfauthor={Autorenteam},
  pdfsubject={Interdisziplinäre Projektarbeit 2026},
  pdflang={de-CH}
]{hyperref}
\usepackage{cleveref}             % Intelligente Querverweise (\cref)

% -----------------------------------------------------------------------------
% Literaturverzeichnis (BibLaTeX mit biber)
% -----------------------------------------------------------------------------
\usepackage[
  backend=biber,
  style=ieee,
  sorting=none
]{biblatex}
\addbibresource{literatur.bib}    % BibTeX-Datenbank

% -----------------------------------------------------------------------------
% Kopf- und Fusszeile
% -----------------------------------------------------------------------------
\usepackage{scrlayer-scrpage}
\pagestyle{scrheadings}
\clearpairofpagestyles
\automark[section]{chapter}
\ihead{\headmark}                 % Innen: aktuelles Kapitel
\ohead{\pagemark}                 % Aussen: Seitenzahl
\cfoot{}

% =============================================================================
% Dokument beginnt hier
% =============================================================================
\begin{document}

% --- Titelblatt (keine Seitenzahl) -------------------------------------------
% =============================================================================
% 00_titelblatt.tex — Titelblatt der IDPA
% =============================================================================
% Dieses Blatt bildet die offizielle Titelseite gemäss Schulvorgabe.
% Ersetze alle Platzhalter (<<...>>) durch die tatsächlichen Angaben.
% =============================================================================

\begin{titlepage}
  \centering

  % --- Logo und Name der Schule ----------------------------------------------
  % Binde das offizielle Schullogo ein (Datei z. B. images/schullogo.png).
  % Falls kein Logo vorhanden ist, kann die Zeile \includegraphics auskommentiert
  % und durch den Schulnamen in Textform ersetzt werden.
  % Passe die Dateiendung an das tatsächliche Format an (.png, .pdf, .jpg)
  \includegraphics[width=5cm]{schullogo.png}\\[0.5cm]
  {\Large \textbf{<<Name der Schule>>}}\\[0.3cm]
  {\large <<Abteilung / Fachrichtung>>}

  \vspace{2cm}

  % --- Art der Arbeit --------------------------------------------------------
  {\large Interdisziplinäre Projektarbeit (IDPA)}\\[0.5cm]

  % --- Vollständiger Titel der Arbeit ----------------------------------------
  {\Huge \textbf{Geotracker}}\\[0.4cm]
  {\Large \textbf{<<Untertitel der Arbeit>>}}

  \vspace{2cm}

  % --- Autor:innen und Klasse ------------------------------------------------
  % Füge alle beteiligten Personen mit Klasse ein.
  \begin{tabular}{ll}
    \textbf{Autor:innen:} & <<Vorname Nachname>>, Klasse <<Klasse>> \\
                          & <<Vorname Nachname>>, Klasse <<Klasse>> \\
                          & <<Vorname Nachname>>, Klasse <<Klasse>> \\
  \end{tabular}

  \vspace{1.5cm}

  % --- Betreuende Lehrpersonen -----------------------------------------------
  \begin{tabular}{ll}
    \textbf{Betreuung:} & <<Titel Vorname Nachname>> \\
                        & <<Titel Vorname Nachname>> \\
  \end{tabular}

  \vspace{1.5cm}

  % --- Abgabedatum -----------------------------------------------------------
  \begin{tabular}{ll}
    \textbf{Abgabedatum:} & <<TT. Monat JJJJ>> \\
  \end{tabular}

  \vfill

  % --- Schuljahr -------------------------------------------------------------
  {\small Schuljahr <<JJJJ/JJJJ>>}

\end{titlepage}

% Seitenzähler zurücksetzen, damit das Titelblatt keine Seitenzahl trägt
\setcounter{page}{0}


% --- Frontmatter (römische Seitenzahlen) -------------------------------------
\frontmatter
% =============================================================================
% 01_abstract.tex — Abstract / Zusammenfassung
% =============================================================================
% Prägnante Zusammenfassung der gesamten Arbeit (max. eine halbe A4-Seite).
% Inhalt: Was wurde analysiert? Welches Vorgehen wurde gewählt?
%         Welche Erkenntnisse wurden gewonnen?
% Der Abstract richtet sich an Lesende, die die Arbeit noch nicht kennen.
% =============================================================================

\chapter*{Abstract}
\addcontentsline{toc}{chapter}{Abstract}  % Manuell ins Inhaltsverzeichnis einfügen

% --- Analysiertes Problem / Ausgangslage ------------------------------------
% Beschreibe in 2–3 Sätzen, welches Problem oder welche Fragestellung
% den Ausgangspunkt der Arbeit bildet.
<<Hier die Ausgangslage und das analysierte Problem beschreiben.>>

% --- Gewähltes Vorgehen ------------------------------------------------------
% Erkläre kurz, wie das Problem angegangen wurde (Methoden, Technologien).
<<Hier das methodische Vorgehen und die eingesetzten Technologien beschreiben.>>

% --- Erzielte Erkenntnisse / Ergebnisse --------------------------------------
% Fasse die wichtigsten Resultate und Erkenntnisse der Arbeit zusammen.
<<Hier die wichtigsten Ergebnisse und Erkenntnisse zusammenfassen.>>

\tableofcontents
\listoffigures
\listoftables

% --- Hauptteil (arabische Seitenzahlen) --------------------------------------
\mainmatter
% =============================================================================
% 02_einleitung.tex — Einleitung
% =============================================================================
% Die Einleitung führt die Lesenden in die Arbeit ein. Sie enthält:
%   1. Grundlagen und Definition des Themas
%   2. Ziele der Arbeit
%   3. Technisches Vorgehen / Anforderungsspezifikation
%   4. Aufbau der Arbeit
% =============================================================================

\chapter{Einleitung}

% --- 1. Grundlagen und Definition des Themas ---------------------------------
% Erkläre den fachlichen Kontext. Was ist ein Geotracker?
% Welche technologischen Grundlagen (GPS, IoT, Cloud, Web) sind relevant?
% Definiere zentrale Begriffe, die in der Arbeit verwendet werden.
\section{Grundlagen und Themendefinition}

<<Hier den technischen und fachlichen Hintergrund des Geotrackers beschreiben.
Zentrale Begriffe definieren (z.\,B.\ GPS, GNSS, MQTT, REST-API, Progressive
Web App).>>

% --- 2. Ziele der Arbeit -----------------------------------------------------
% Formuliere klare, messbare Ziele (z. B. als Forschungsfragen oder Muss-Ziele).
\section{Ziele der Arbeit}

<<Hier die konkreten Ziele und Forschungsfragen der IDPA formulieren.
Beispiel: Kann ein kostengünstiger GPS-Tracker mit Echtzeit-Webdarstellung
realisiert werden?>>

% --- 3. Technisches Vorgehen / Anforderungsspezifikation ---------------------
% Beschreibe den gewählten Entwicklungsansatz (z. B. agil, iterativ).
% Liste funktionale und nicht-funktionale Anforderungen auf.
\section{Technisches Vorgehen und Anforderungsspezifikation}

\subsection{Funktionale Anforderungen}
% Was muss das System können? (Use-Cases, User Stories o. ä.)
<<Funktionale Anforderungen als Liste oder Tabelle einfügen.>>

\subsection{Nicht-funktionale Anforderungen}
% Leistung, Sicherheit, Zuverlässigkeit, Wartbarkeit usw.
<<Nicht-funktionale Anforderungen beschreiben.>>

% --- 4. Aufbau der Arbeit ----------------------------------------------------
% Erkläre kurz, was in den folgenden Kapiteln behandelt wird.
\section{Aufbau der Arbeit}

Die vorliegende Arbeit gliedert sich wie folgt:
\begin{description}
  \item[Kapitel~\ref{chap:hauptteil} – Hauptteil]
    Beschreibt die Entwicklung in den drei technischen Domänen Firmware,
    Cloud-Konfiguration und Web-Frontend.
  \item[Kapitel~\ref{chap:fazit} – Fazit]
    Bewertet die erzielten Ergebnisse kritisch und beantwortet die
    Forschungsfragen.
  \item[Anhang]
    Enthält ergänzende Unterlagen, detaillierte Auswertungen und
    unveröffentlichte Quellen.
\end{description}

% =============================================================================
% 03_hauptteil_main.tex — Hauptteil / Vertiefung
% =============================================================================
% Dieses Kapitel ist das Kernstück der Arbeit. Es ist in drei technische
% Domänen aufgeteilt, die als separate Dateien eingebunden werden:
%   03_1_hauptteil_firmware.tex       — Firmware (Mikrocontroller / Embedded)
%   03_2_hauptteil_cloud_config.tex   — Cloud-Konfiguration & Backend
%   03_3_hauptteil_web_frontend.tex   — Web-Frontend (UI / Dashboard)
% =============================================================================

\chapter{Hauptteil / Vertiefung}
\label{chap:hauptteil}

% Kurze Einführung in den Hauptteil
% Erkläre die Gesamtarchitektur des Systems und wie die drei Domänen
% zusammenwirken. Ein Architekturdiagramm ist hier empfehlenswert.
<<Gesamtarchitektur des Geotrackers kurz beschreiben und Zusammenspiel der
drei Teilsysteme erläutern.>>

% Architekturdiagramm (optional, Platzhalter):
% \begin{figure}[htbp]
%   \centering
%   \includegraphics[width=0.9\textwidth]{architektur_gesamt}
%   \caption{Gesamtarchitektur des Geotracker-Systems}
%   \label{fig:architektur_gesamt}
% \end{figure}

% --- Domäne 1: Firmware -------------------------------------------------------
% =============================================================================
% 03_1_hauptteil_firmware.tex — Domäne: Firmware
% =============================================================================
% Dieser Abschnitt dokumentiert die Embedded-Software (Firmware) des
% Geotracker-Geräts. Gliederung gemäss Schulvorgabe:
%   1. Analyse und Design
%   2. Softwaredokumentation
%   3. Reflexion
% =============================================================================

\section{Firmware}
\label{sec:firmware}

% Kurze Einführung in die Firmware-Domäne
% Welcher Mikrocontroller / welches Modul wird verwendet?
% Welche Hauptaufgabe erfüllt die Firmware?
<<Kurze Einführung in die verwendete Hardware (z.\,B.\ ESP32, u-blox GPS-Modul)
und die Aufgabe der Firmware innerhalb des Gesamtsystems.>>

% --- 1. Analyse und Design ---------------------------------------------------
\subsection{Analyse und Design}

% Beschreibe die Analysephase: Anforderungen an die Firmware,
% Auswahl der Hardware-Komponenten, Entscheidungen zur Architektur.
% Füge UML-Diagramme (Zustandsdiagramm, Komponentendiagramm) ein,
% die den Ablauf und die Struktur der Firmware zeigen.

\subsubsection{Anforderungsanalyse}
<<Anforderungen an die Firmware auflisten (z.\,B.\ GPS-Daten auslesen,
Koordinaten per MQTT übertragen, Energiesparmodus).>>

\subsubsection{Systemarchitektur und UML-Diagramme}
% Platzhalter für UML-Diagramme (Zustandsautomat, Komponentendiagramm):
% \begin{figure}[htbp]
%   \centering
%   \includegraphics[width=0.85\textwidth]{uml_firmware_zustand}
%   \caption{Zustandsdiagramm der Firmware}
%   \label{fig:uml_firmware}
% \end{figure}
<<UML-Diagramme (Zustands- und Komponentendiagramm) hier einfügen und erläutern.>>

% --- 2. Softwaredokumentation ------------------------------------------------
\subsection{Softwaredokumentation}

\subsubsection{Beschreibung und Konzepte}
% Erkläre die wichtigsten Softwaremodule, Bibliotheken und Konzepte.
% Wie werden GPS-Daten verarbeitet? Wie funktioniert die MQTT-Übertragung?
<<Hauptmodule der Firmware beschreiben (z.\,B.\ GPS-Parser, MQTT-Client,
Schlafmodus-Verwaltung). Wichtige Codeausschnitte als Listings einfügen.>>

% Beispiel für ein Code-Listing:
% \begin{lstlisting}[caption={GPS-Daten auslesen}, label={lst:gps_read}]
%   void readGPS() {
%     // GPS-Daten vom seriellen Port lesen
%   }
% \end{lstlisting}

\subsubsection{Schwierigkeiten und Lösungsansätze}
% Welche technischen Probleme traten auf? Wie wurden sie gelöst?
<<Technische Herausforderungen dokumentieren (z.\,B.\ GPS-Signalverlust in
Gebäuden, Timing-Probleme bei MQTT-Verbindung) und die gewählten
Lösungsstrategien beschreiben.>>

\subsubsection{Testkonzept}
% Wie wurde die Firmware getestet? (Unit-Tests, Integrationstests, Feldtests)
% Welche Testfälle wurden definiert? Wie wurden die Ergebnisse dokumentiert?
<<Testkonzept und Testergebnisse für die Firmware dokumentieren.
Tabelle mit Testfällen und deren Status (bestanden / fehlgeschlagen) einfügen.>>

% --- 3. Reflexion ------------------------------------------------------------
\subsection{Reflexion}
% Bewerte die eigene Arbeit an der Firmware kritisch.
% Was lief gut? Was würde man anders machen?
% Inwiefern wurden die Anforderungen erfüllt?
<<Persönliche und technische Reflexion zur Firmware-Entwicklung.
Stärken und Schwächen des gewählten Ansatzes benennen.>>


% --- Domäne 2: Cloud-Konfiguration --------------------------------------------
% =============================================================================
% 03_2_hauptteil_cloud_config.tex — Domäne: Cloud-Konfiguration
% =============================================================================
% Dieser Abschnitt dokumentiert die Cloud-Infrastruktur und Backend-
% Konfiguration des Geotracker-Projekts. Gliederung gemäss Schulvorgabe:
%   1. Analyse und Design
%   2. Softwaredokumentation
%   3. Reflexion
% =============================================================================

\section{Cloud-Konfiguration}
\label{sec:cloud}

% Kurze Einführung in die Cloud-Domäne
% Welche Cloud-Plattform wird genutzt? Welche Dienste sind im Einsatz?
% (z. B. MQTT-Broker, Datenbank, API-Server)
<<Kurze Einführung in die Cloud-Architektur (z.\,B.\ AWS IoT Core, Azure IoT Hub,
selbst gehosteter MQTT-Broker, Datenbank) und deren Rolle im Gesamtsystem.>>

% --- 1. Analyse und Design ---------------------------------------------------
\subsection{Analyse und Design}

\subsubsection{Anforderungsanalyse}
% Welche Anforderungen hat das Cloud-Backend zu erfüllen?
% (Skalierbarkeit, Verfügbarkeit, Datenschutz, Kosten)
<<Anforderungen an das Cloud-Backend auflisten (z.\,B.\ Empfang von GPS-Daten
per MQTT, Persistierung in Datenbank, REST-API für das Web-Frontend,
Zugriffskontrolle).>>

\subsubsection{Architekturdesign und Diagramme}
% Skizziere die Cloud-Architektur (Komponenten, Datenfluss).
% Füge ein Architektur- oder Sequenzdiagramm ein.
% \begin{figure}[htbp]
%   \centering
%   \includegraphics[width=0.85\textwidth]{uml_cloud_architektur}
%   \caption{Architekturdiagramm der Cloud-Konfiguration}
%   \label{fig:cloud_architektur}
% \end{figure}
<<Architektur- und Sequenzdiagramme hier einfügen und den Datenfluss von der
Firmware bis zur Web-Anwendung erläutern.>>

% --- 2. Softwaredokumentation ------------------------------------------------
\subsection{Softwaredokumentation}

\subsubsection{Beschreibung und Konzepte}
% Dokumentiere die eingesetzten Dienste, Protokolle und Konfigurationen.
% Wie ist der MQTT-Broker konfiguriert? Welche Datenbankstruktur wird verwendet?
<<Dienste und Konfigurationen beschreiben (z.\,B.\ MQTT-Topics, Datenbankschema,
API-Endpunkte, Authentifizierungsmechanismus). Relevante Konfigurationsdateien
oder Codeausschnitte als Listings einfügen.>>

% Beispiel:
% \begin{lstlisting}[caption={MQTT-Topic-Struktur}, label={lst:mqtt_topics},
%                    language={}]
%   geotracker/<device-id>/location   % GPS-Koordinaten
%   geotracker/<device-id>/status     % Gerätestatus
% \end{lstlisting}

\subsubsection{Schwierigkeiten und Lösungsansätze}
% Welche Herausforderungen traten bei der Cloud-Konfiguration auf?
% (z. B. Latenz, Authentifizierung, Skalierung, Kosten)
<<Technische Herausforderungen dokumentieren (z.\,B.\ TLS-Zertifikatsverwaltung,
Datenbank-Performance bei vielen GPS-Punkten) und Lösungsstrategien beschreiben.>>

\subsubsection{Testkonzept}
% Wie wurde das Cloud-Backend getestet?
% (End-to-End-Tests, Lasttests, Sicherheitstests)
<<Testkonzept für das Cloud-Backend beschreiben. Testergebnisse und
verwendete Tools dokumentieren (z.\,B.\ MQTT-Explorer, Postman für API-Tests).>>

% --- 3. Reflexion ------------------------------------------------------------
\subsection{Reflexion}
% Bewerte die Cloud-Lösung kritisch.
% Würdest du dieselbe Plattform wieder wählen? Was würdest du verbessern?
<<Reflexion zur Cloud-Konfiguration: Stärken (z.\,B.\ gute Skalierbarkeit),
Schwächen (z.\,B.\ hohe Komplexität, Kosten) und Verbesserungsvorschläge.>>


% --- Domäne 3: Web-Frontend ---------------------------------------------------
% =============================================================================
% 03_3_hauptteil_web_frontend.tex — Domäne: Web-Frontend
% =============================================================================
% Dieser Abschnitt dokumentiert das Web-Frontend (Dashboard / UI) des
% Geotracker-Projekts. Gliederung gemäss Schulvorgabe:
%   1. Analyse und Design
%   2. Softwaredokumentation
%   3. Reflexion
% =============================================================================

\section{Web-Frontend}
\label{sec:frontend}

% Kurze Einführung in das Web-Frontend
% Welches Framework / welche Technologien werden eingesetzt?
% Was zeigt das Frontend dem Benutzer?
<<Kurze Einführung in das Web-Frontend (z.\,B.\ React, Vue.js, Leaflet.js für
Kartenansicht) und dessen Funktion im Gesamtsystem (Echtzeit-Positionsanzeige,
Verlaufsdarstellung).>>

% --- 1. Analyse und Design ---------------------------------------------------
\subsection{Analyse und Design}

\subsubsection{Anforderungsanalyse}
% Welche Anforderungen hat das Frontend zu erfüllen?
% (Benutzerfreundlichkeit, Responsivität, Echtzeit-Updates)
<<Anforderungen an das Web-Frontend auflisten (z.\,B.\ Kartenansicht mit
Echtzeit-Position, Verlaufslinie, Filterfunktionen, mobile Kompatibilität).>>

\subsubsection{GUI-Skizzen und Mockups}
% Füge Skizzen oder Mockups der wichtigsten Ansichten ein.
% Zeige, wie die Benutzeroberfläche aufgebaut ist.
% \begin{figure}[htbp]
%   \centering
%   \includegraphics[width=0.85\textwidth]{mockup_dashboard}
%   \caption{GUI-Mockup des Geotracker-Dashboards}
%   \label{fig:mockup_dashboard}
% \end{figure}
<<GUI-Skizzen / Wireframes für die Hauptansichten (Dashboard, Detailansicht,
Einstellungen) hier einfügen und erläutern.>>

\subsubsection{Komponentendiagramm}
% Skizziere die Komponentenstruktur der Frontend-Anwendung (UML oder informell).
<<Komponentendiagramm des Frontends beschreiben: Welche Komponenten gibt es?
Wie kommunizieren sie miteinander und mit der API?>>

% --- 2. Softwaredokumentation ------------------------------------------------
\subsection{Softwaredokumentation}

\subsubsection{Beschreibung und Konzepte}
% Erkläre die wichtigsten Komponenten, Zustandsverwaltung, Datenabruf.
% Wie werden Echtzeit-Daten abgerufen (WebSocket, SSE, Polling)?
<<Hauptkomponenten des Frontends beschreiben (z.\,B.\ Kartenkomponente,
Positionsmarker, Datenabruf via WebSocket/REST). Wichtige Codeausschnitte
als Listings einfügen.>>

% Beispiel:
% \begin{lstlisting}[caption={WebSocket-Verbindung aufbauen},
%                    label={lst:websocket}, language=JavaScript]
%   const ws = new WebSocket('wss://api.example.com/live');
%   ws.onmessage = (event) => {
%     const { lat, lng } = JSON.parse(event.data);
%     updateMarker(lat, lng);
%   };
% \end{lstlisting}

\subsubsection{Schwierigkeiten und Lösungsansätze}
% Welche Herausforderungen traten bei der Frontend-Entwicklung auf?
% (z. B. Cross-Origin-Requests, Kartenperformance, mobile Darstellung)
<<Technische Herausforderungen dokumentieren (z.\,B.\ CORS-Probleme, Performance
bei vielen GPS-Punkten auf der Karte, Responsive Design) und Lösungsstrategien
beschreiben.>>

\subsubsection{Testkonzept}
% Wie wurde das Frontend getestet?
% (Manuelle Tests, automatisierte UI-Tests, Cross-Browser-Tests)
<<Testkonzept für das Web-Frontend beschreiben. Welche Browser und Geräte wurden
getestet? Welche automatisierten Tests wurden eingesetzt (z.\,B.\ Jest,
Playwright)?>>

% --- 3. Reflexion ------------------------------------------------------------
\subsection{Reflexion}
% Bewerte das Frontend kritisch.
% Ist die Benutzeroberfläche intuitiv? Was würdest du verbessern?
<<Reflexion zur Frontend-Entwicklung: Stärken (z.\,B.\ übersichtliche Karte),
Schwächen (z.\,B.\ fehlende Offline-Unterstützung) und Verbesserungsvorschläge.>>


% =============================================================================
% 04_fazit.tex — Fazit und Schlussbetrachtung
% =============================================================================
% Das Fazit bewertet die Arbeit kritisch und abschliessend. Inhalt:
%   1. Kritische Überprüfung der Resultate und des Prozesses
%   2. Beantwortung der Forschungsfragen
%   3. Lösungsstrategien (Rückblick)
%   4. Ergebnisdokumentation / Ausblick
% =============================================================================

\chapter{Fazit}
\label{chap:fazit}

% --- 1. Kritische Überprüfung der Resultate ----------------------------------
\section{Kritische Überprüfung der Resultate}

% Wurden die in der Einleitung gesetzten Ziele erreicht?
% Wo liegen die Stärken, wo die Grenzen des realisierten Systems?
% Vergleiche die erzielten Resultate mit dem ursprünglichen Soll.
<<Hier die erreichten Ergebnisse mit den definierten Zielen aus der Einleitung
vergleichen. Was funktioniert gut? Welche Anforderungen wurden nur teilweise
oder gar nicht erfüllt?>>

% --- 2. Beantwortung der Forschungsfragen ------------------------------------
\section{Beantwortung der Forschungsfragen}

% Beantworte jede in der Einleitung formulierte Forschungsfrage einzeln
% und direkt, gestützt auf die Erkenntnisse aus dem Hauptteil.
<<Forschungsfragen aus der Einleitung einzeln aufgreifen und anhand der
erarbeiteten Ergebnisse beantworten. Belege die Antworten mit konkreten
Daten oder Beobachtungen.>>

% --- 3. Lösungsstrategien (Rückblick) ----------------------------------------
\section{Rückblick auf die Lösungsstrategien}

% Reflektiere, ob die gewählten Ansätze (technologische Entscheide,
% Entwicklungsmethode, Aufteilung der Aufgaben) die richtigen waren.
% Was würde das Team beim nächsten Mal anders machen?
<<Den Entwicklungsprozess und die gewählten Lösungsstrategien rückblickend
bewerten. Welche Entscheidungen erwiesen sich als richtig, welche als
problematisch?>>

% --- 4. Ausblick / Ergebnisdokumentation -------------------------------------
\section{Ausblick}

% Was sind die nächsten sinnvollen Schritte, falls das Projekt weitergeführt
% würde? Welche Features oder Verbesserungen wären erstrebenswert?
<<Mögliche Weiterentwicklungen und offene Punkte benennen (z.\,B.\ Erweiterung
um Geofencing, Optimierung des Energieverbrauchs, Skalierung des Backends).>>

% =============================================================================
% 05_ehrenwoertliche_erklaerung.tex — Ehrenwörtliche Erklärung
% =============================================================================
% Dieses Kapitel enthält die Bestätigung der selbständigen Erarbeitung.
% WICHTIG: Gemäss Schulvorgabe muss explizit erwähnt werden, dass die
% Bestimmungen zur Selbständigkeit auch für den Einsatz von
% Künstlicher Intelligenz (KI) als Hilfsmittel gelten.
% =============================================================================

\chapter*{Ehrenwörtliche Erklärung}
\addcontentsline{toc}{chapter}{Ehrenwörtliche Erklärung}

Wir erklären hiermit ehrenwörtlich, dass wir die vorliegende
Interdisziplinäre Projektarbeit (IDPA) selbständig und ohne unerlaubte
fremde Hilfe verfasst haben. Alle verwendeten Quellen und Hilfsmittel
sind vollständig angegeben und im Literaturverzeichnis aufgeführt.

Wörtliche oder sinngemässe Übernahmen aus anderen Werken sind als solche
kenntlich gemacht. Die Arbeit wurde weder in gleicher noch in ähnlicher
Form einer anderen Prüfungsbehörde vorgelegt.

% --- Hinweis zu Künstlicher Intelligenz (gemäss Schulvorgabe) ----------------
% WICHTIG: Dieser Absatz ist gemäss den Bestimmungen der Schule obligatorisch.
Uns ist bekannt, dass die Bestimmungen zur ehrenwörtlichen Erklärung und
zur Eigenständigkeit der Arbeit auch für den Einsatz von Hilfsmitteln in
Form von \textbf{Künstlicher Intelligenz (KI)} gelten. Der Einsatz von
KI-gestützten Werkzeugen (z.\,B.\ Large-Language-Models, KI-basierte
Code-Assistenten oder Text\-generierungs\-dienste) ist gemäss den
geltenden Richtlinien der Schule zu deklarieren und entspricht einer
fremden Hilfe, sofern die daraus resultierenden Inhalte ohne
eigenständige Reflexion und kritische Überarbeitung übernommen werden.

Allfällig eingesetzte KI-Hilfsmittel sowie der Umfang ihrer Verwendung
sind nachfolgend vollständig offengelegt:

\begin{itemize}
  % Füge hier alle verwendeten KI-Werkzeuge ein oder trage «keine» ein.
  \item <<Name des KI-Werkzeugs (z.\,B.\ GitHub Copilot, ChatGPT)>>,
        eingesetzt für: <<Beschreibung des Verwendungszwecks und -umfangs>>
  \item <<weiteres KI-Werkzeug oder «keine weiteren KI-Hilfsmittel»>>
\end{itemize}

\vspace{2cm}

% Unterschriftenfelder — Ort, Datum und Unterschriften aller Autor:innen
\noindent
Ort, Datum: \underline{\hspace{6cm}}

\vspace{1.5cm}
\noindent
\begin{tabular}{p{6cm} p{6cm}}
  \underline{\hspace{5.5cm}} & \underline{\hspace{5.5cm}} \\[0.3cm]
  <<Vorname Nachname>>        & <<Vorname Nachname>>        \\[1.5cm]
  \underline{\hspace{5.5cm}} &                             \\[0.3cm]
  <<Vorname Nachname>>        &                             \\
\end{tabular}

% =============================================================================
% 06_anhang.tex — Anhang
% =============================================================================
% Der Anhang enthält ergänzende Materialien, die den Lesefluss im Hauptteil
% gestört hätten. Typische Inhalte:
%   - Selbst erstellte Unterlagen (Schaltpläne, Datenblätter, Protokolle)
%   - Unveröffentlichte Quellen (interne Dokumente, E-Mails, Interviews)
%   - Detaillierte Auswertungen (Rohdaten, umfangreiche Tabellen)
%   - Vollständige Quelltexte (sofern nicht im Repo verlinkt)
% =============================================================================

% Anhang-Formatierung: Kapitel werden mit Buchstaben nummeriert (A, B, C …)
\appendix

\chapter{Anhang}
\label{chap:anhang}

% --- A. Schaltpläne und Hardware-Dokumentation --------------------------------
\section{Schaltpläne und Hardware-Dokumentation}
% Füge Schaltpläne, Stücklisten (BOM) oder Datenblätter ein.
% \begin{figure}[htbp]
%   \centering
%   \includegraphics[width=0.9\textwidth]{schaltplan_geotracker}
%   \caption{Schaltplan des Geotracker-Prototyps}
%   \label{fig:schaltplan}
% \end{figure}
<<Schaltpläne, Stücklisten (BOM) und weitere Hardware-Unterlagen hier einfügen.>>

% --- B. Messprotokolle und Testergebnisse ------------------------------------
\section{Messprotokolle und Testergebnisse}
% Detaillierte Tabellen mit Testfällen und Messwerten, die im Hauptteil
% zu umfangreich wären.
<<Vollständige Testprotokolle, Messwert-Tabellen und Auswertungen hier einfügen.>>

% --- C. Unveröffentlichte Quellen --------------------------------------------
\section{Unveröffentlichte Quellen}
% Interne Dokumente, E-Mails, Interview-Transkripte, die nicht öffentlich
% zugänglich sind, aber als Quellen dienen.
<<Unveröffentlichte Quellen (mit Datum und Kontext) hier einfügen.>>

% --- D. Weitere Unterlagen ---------------------------------------------------
\section{Weitere Unterlagen}
% Platzhalter für sonstige relevante Materialien.
<<Weitere ergänzende Unterlagen hier einfügen.>>


% --- Backmatter --------------------------------------------------------------
\backmatter
\printbibliography[heading=bibintoc, title={Literaturverzeichnis}]

\end{document}
